\newcolumntype{C}[1]{>{\centering}p{#1}}
\newcolumntype{L}[1]{>{\raggedleft}p{#1}}
\newcolumntype{C}[1]{>{\centering}p{#1}}
\newcolumntype{L}[1]{>{\raggedleft}p{#1}}
\small
\begin{tabularx}{\textwidth}{L{\firstColWidth{}}|C{\secondColWidth{}}|X}
\toprule
{\bf RPM Package Name} & {\bf Version} & {\bf Info/URL}  \\
\midrule

% <-- begin entry for gnu15-compilers-ohpc
\multirow{2}{*}{gnu15-compilers-ohpc} &
\multirow{2}{*}{15.2.0} &
GNU Compiler Collection 15.2.0. \newline { \color{logoblue} \url{https://gcc.gnu.org}}
\\ \hline
% <-- end entry for gnu15-compilers-ohpc

% <-- begin entry for intel-compilers-devel-ohpc
\multirow{2}{*}{intel-compilers-devel-ohpc} &
\multirow{2}{*}{2025.0} &
OpenHPC compatibility package for Intel(R) oneAPI HPC Toolkit.  { \color{logoblue} \url{https://github.com/openhpc/ohpc}}
\\ \hline
% <-- end entry for intel-compilers-devel-ohpc

% <-- begin entry for intel-oneapi-toolkit-release-ohpc
\multirow{2}{*}{intel-oneapi-toolkit-release-ohpc} &
\multirow{2}{*}{2025.0} &
Intel(R) oneAPI HPC Toolkit Repository Setup. \newline { \color{logoblue} \url{https://github.com/openhpc/ohpc}}
\\ \hline
% <-- end entry for intel-oneapi-toolkit-release-ohpc

% <-- begin entry for llvm10-compilers-ohpc
\multirow{2}{*}{llvm10-compilers-ohpc} &
\multirow{2}{*}{10.0.1} &
LLVM compiler infrastructure for OpenHPC. \newline { \color{logoblue} \url{https://llvm.org}}
\\ \hline
% <-- end entry for llvm10-compilers-ohpc

\bottomrule
\end{tabularx}
